\problemname{Miniräknaren}
% Statement is kind of todo
Nyligen har det gått runt rykten på skolan att skolans miniräknare kan trimmas med en hemlig kod. 
Det sägs att om du kan få miniräknaren att visa den hemliga kålnami-koden på displayen och sedan
trycka på \texttt{=} tecknet, så kommer miniräknaren att låsa upp den förbjudna kvadratrotsfunktionen.
Detta kompliceras ytterligare av att du måste skriva in den på ett särskilt sett: du börjar med att
miniräknaren visar $0$. Sedan får du utföra följande operationer:
\begin{itemize}
  \item Gångra talet som visas med siffran $m$
  \item Addera en valfri siffra $k$ till talet, där $k$ uppfyller $0 \leq k \leq m-1$
\end{itemize}

För att göra det intressant funkar det bara om du gör detta med så få operationer som möjligt. Din kompis
har precis hittat kålnami-koden och siffran $m$, nu är det bara att hitta minsta antalet operationer
som krävs.

\section*{Indata}
Den första raden innehåller heltalet $n$ ($1 \le n \le 10^9$), innehållet av kolnami-koden.

Nästa rad innehåller siffran $m$ ($2 \leq m \leq 9$), som beskrivits ovan.

\section*{Utdata}
Skriv ut ett heltal: minsta antalet operationer som krävs för att få miniräknaren att visa
kålnami-koden i sina siffer-display.


\section*{Poängsättning}
Din lösning kommer att testas på en mängd testfallsgrupper.
För att få poäng för en grupp så måste du klara alla testfall i gruppen.

\noindent
\begin{tabular}{| l | l | p{12cm} |}
  \hline
  \textbf{Grupp} & \textbf{Poäng} & \textbf{Gränser} \\ \hline
  $1$    & $20$       & $n \leq 10, m=3$. \\ \hline
  $2$    & $20$       & $m=2$ \\ \hline
  $3$    & $40$       & $n \leq 10^5$ \\ \hline
  $4$    & $20$       & Inga ytterligare begränsningar. \\ \hline
\end{tabular}
